\subsection{The unanchored dependence and the anchored kernel family}\label{sec:binary:anchors}

The identity $\sum_bK_b=1$ creates an exact linear
dependence between the constant polynomial and the
kernel basis. This makes a naive residual polynomial
identity test unsound unless the additional scalar
residual is distinguished algebraically.
We derive the minimal-degree anchoring mechanism;
a full protocol using it will appear later.

\begin{lemma}[Unanchored residual dependence]\label{lem:binary:unanchored}
Let $\delta\in L$ be nonzero and set $e_b:=\delta$ for
every $b\in I_\ell$. Then
\[
 \delta+\sum_{b\in I_\ell}e_bK_b(X)=0
\]
as a polynomial, even though $(\delta,e)\neq0$.
\end{lemma}
\begin{proof}
By Lemma~\ref{lem:binary:partition},
$\sum_bK_b=1$. Therefore the left side is
$\delta+\delta\sum_bK_b=\delta+\delta=0$
in characteristic two. The first residual component
$\delta$ is nonzero by assumption.
\end{proof}

\begin{definition}[Anchored Boolean kernels]\label{def:binary:anchored}
Let $0_\ell=(0,\ldots,0)\in I_\ell$. For $b\in I_\ell$ define
\[
 W_b(X):=
 \begin{cases}
 K_{0_\ell}(X)+X^n,&b=0_\ell,\\
 K_b(X),&b\neq0_\ell.
 \end{cases}
\]
Every $W_b$ is a polynomial with coefficients in $\F_2$.
\end{definition}

\begin{lemma}[Anchored family degree]\label{lem:binary:anchor-degree}
$\deg W_{0_\ell}=n$, while
$\deg W_b=n-1$ for every $b\neq0_\ell$.
\end{lemma}
\begin{proof}
By Lemma~\ref{lem:intro:degree}, each kernel has degree
$n-1$. Therefore the added monomial $X^n$ is the
unique degree-$n$ contribution to $W_{0_\ell}$,
with nonzero coefficient $1$, so its degree is $n$.
For $b\neq0_\ell$, $W_b=K_b$, which has degree $n-1$.
\end{proof}

\begin{lemma}[Sum of anchored kernels]\label{lem:binary:anchor-sum}
\[
 \sum_{b\in I_\ell}W_b(X)=1+X^n.
\]
\end{lemma}
\begin{proof}
By definition only the term $W_{0_\ell}$
differs from $K_{0_\ell}$, by the summand $X^n$.
Thus $\sum_bW_b=\sum_bK_b+X^n=1+X^n$
by Lemma~\ref{lem:binary:partition}.
\end{proof}

\begin{theorem}[Independence of the constant and all anchored kernels]\label{thm:binary:anchor-independence}
The $n+1$ polynomials
\[
 \boxed{\{1\}\cup\{W_b:b\in I_\ell\}}
\]
are linearly independent over $L$ in $L[X]_{\leq n}$.
\end{theorem}
\begin{proof}
Suppose
\[
 \delta+\sum_{b\in I_\ell}e_bW_b(X)=0
 \qquad(\delta,e_b\in L).
\]
The coefficient of $X^n$ in the left side is
exactly $e_{0_\ell}$ by Lemma~\ref{lem:binary:anchor-degree}.
Hence $e_{0_\ell}=0$. The remaining polynomial identity is
$\delta+\sum_{b\neq0_\ell}e_bK_b=0$.
Using $1=\sum_bK_b$ from
Lemma~\ref{lem:binary:partition}, rewrite it as
\[
 \delta K_{0_\ell}
 +\sum_{b\neq0_\ell}(\delta+e_b)K_b=0.
\]
By the kernel basis theorem,
Proposition~\ref{prop:intro:basis},
all kernel coefficients vanish. In particular $\delta=0$,
and then $\delta+e_b=0$ gives every $e_b=0$.
Thus the anchored family together with $1$ is independent.
\end{proof}

\begin{corollary}[Anchored family is a basis in degree at most $n$]\label{cor:binary:anchor-basis}
The $n+1$ polynomials from
Theorem~\ref{thm:binary:anchor-independence}
form an $L$-basis of $L[X]_{\leq n}$.
\end{corollary}
\begin{proof}
The monomials $1,X,\ldots,X^n$ form a basis of
$L[X]_{\leq n}$, which therefore has dimension $n+1$.
The displayed family consists of $n+1$ linearly
independent vectors in this space, hence is a basis.
\end{proof}

\begin{proposition}[Degree-$n$ anchoring is dimensionally minimal]\label{prop:binary:anchor-minimal}
No family of $n+1$ polynomials, all of degree $<n$,
can be linearly independent over $L$.
\end{proposition}
\begin{proof}
The space $L[X]_{<n}$ has basis
$1,X,\ldots,X^{n-1}$ and therefore dimension $n$.
A vector space of dimension $n$ cannot contain $n+1$
linearly independent vectors.
\end{proof}

\begin{lemma}[Anchored kernels commute with Frobenius evaluation]\label{lem:binary:anchor-frob}
For every $b\in I_\ell$ and $z\in L$,
\[
 W_b(z)^2=W_b(z^2).
\]
\end{lemma}
\begin{proof}
Every coefficient of $W_b$ lies in $\F_2$
by Definition~\ref{def:binary:anchored}.
The coefficient Frobenius of $W_b$ is therefore $W_b$.
Apply Lemma~\ref{lem:binary:poly-frob} with $j=1$
and evaluate at $z$:
$W_b(z)^2=W_b^{\sigma_1}(z^2)=W_b(z^2)$.
\end{proof}

\begin{definition}[Anchored residual polynomial]\label{def:binary:anchor-residual}
For $\delta\in L$ and $e=(e_b)_{b\in I_\ell}\in L^{I_\ell}$,
define
\[
 F_{\delta,e}(X):=\delta+\sum_{b\in I_\ell}e_bW_b(X).
\]
\end{definition}

\begin{lemma}[Nonzero residual gives a nonzero polynomial]\label{lem:binary:anchor-nonzero}
If $(\delta,e)\neq(0,0)$, then
$F_{\delta,e}$ is a nonzero polynomial of degree at most $n$.
\end{lemma}
\begin{proof}
Each $W_b$ has degree at most $n$,
so the linear combination has degree at most $n$.
If it vanished identically, the linear independence
of $\{1,W_b\}_b$ proved in
Theorem~\ref{thm:binary:anchor-independence}
would force $\delta=0$ and $e_b=0$ for all $b$,
contradicting the hypothesis.
\end{proof}

\begin{theorem}[Root bound for an anchored residual]\label{thm:binary:anchor-root}
Fix $(\delta,e)\neq(0,0)$ independently of a uniform
challenge $z\overset{\$}{\leftarrow}L$. Then
\[
 \boxed{\Pr_z\left[
 \delta+\sum_be_bW_b(z)^2=0\right]\leq\frac{n}{q}.}
\]
\end{theorem}
\begin{proof}
By Lemma~\ref{lem:binary:anchor-frob},
the tested expression is $F_{\delta,e}(z^2)$.
By Corollary~\ref{cor:binary:frob-bijective},
the map $z\mapsto z^2$ is a bijection of the finite
field $L$, so $z^2$ is uniform on $L$.
By Lemma~\ref{lem:binary:anchor-nonzero},
$F_{\delta,e}$ is nonzero of degree at most $n$.
The univariate root bound of
Lemma~\ref{lem:history:rootcount} gives at most $n$
roots. Divide by $q$.
\end{proof}

\begin{remark}[Challenge timing is necessary]
The bound is conditional on $(\delta,e)$ being
fixed independently of the fresh challenge $z$.
If the coefficients are chosen after seeing $z$,
the prover could select a nonzero polynomial
that vanishes at that particular point.
Thus the commit-then-challenge order is part
of the soundness hypothesis, not a cosmetic
presentation detail.
\end{remark}
